\chapter{Lições aprendidas}

\textcolor{red}{Documentar as lições aprendidas de modo a aperfeiçoar os processos e evitar que os erros e problemas encontrados se repitam em futuros projetos.}

\section{Planejado x realizado}
\subsection{Objetivos}
\textcolor{red}{Os objetivos foram alcançados? Responda as questões e comente os pontos mais relevantes – até 700 caracteres, considerando os espaços.}
\subsection{Prazos}
\textcolor{red}{O projeto foi entregue dentro do prazo? Responda as questões e comente os pontos mais relevantes – até 700 caracteres, considerando os espaços.}
\subsection{Orçamento}
\textcolor{red}{O projeto foi entregue dentro do orçamento? Responda as questões e comente os pontos mais relevantes – até 700 caracteres, considerando os espaços.}
\subsection{Escopo}
\textcolor{red}{O  projeto atendeu ao escopo? Responda as questões e comente os pontos mais relevantes – até 700 caracteres, considerando os espaços.}

\section{Projeto}
\textcolor{red}{Comente os pontos mais relevantes a serem aperfeiçoados ou adotados em próximos projetos.}
\subsection{Pontos fortes}
\begin{enumerate}
	\item \textcolor{red}{Até 200 caracteres, considerando os espaços.}
	\item \textcolor{red}{Até 200 caracteres, considerando os espaços.}
	\item \textcolor{red}{Até 200 caracteres, considerando os espaços.}
	\item \textcolor{red}{Até 200 caracteres, considerando os espaços.}
\end{enumerate}
\subsection{Pontos fracos}
\begin{enumerate}
	\item \textcolor{red}{Até 200 caracteres, considerando os espaços.}
	\item \textcolor{red}{Até 200 caracteres, considerando os espaços.}
	\item \textcolor{red}{Até 200 caracteres, considerando os espaços.}
	\item \textcolor{red}{Até 200 caracteres, considerando os espaços.}
\end{enumerate}

\section{Recomendações para projetos futuros}

\begin{enumerate}
	\item \textcolor{red}{Até 200 caracteres, considerando os espaços.}
	\item \textcolor{red}{Até 200 caracteres, considerando os espaços.}
	\item \textcolor{red}{Até 200 caracteres, considerando os espaços.}
	\item \textcolor{red}{Até 200 caracteres, considerando os espaços.}
\end{enumerate}

\section{Questões em aberto}

\textcolor{red}{Usar caso haja alguma questão pendente em relação às entregas do projeto (Ex.: Requisitos não entregues, melhoria na programação do software, troca de sensores, entre outros.)}

\begin{enumerate}
	\item \textcolor{red}{Até 200 caracteres, considerando os espaços.}
	\item \textcolor{red}{Até 200 caracteres, considerando os espaços.}
	\item \textcolor{red}{Até 200 caracteres, considerando os espaços.}
	\item \textcolor{red}{Até 200 caracteres, considerando os espaços.}
\end{enumerate}

\section{Desempenho dos fornecedores}

\subsection{Fornecedores com desempenho acima do esperado}
\textcolor{red}{Fornecedores que você certamente convidaria para um próximo projeto. Citar AAAA – justificativa. Até 200 caracteres, considerando os espaços.}

\subsection{Fornecedores com desempenho conforme esperado}
\textcolor{red}{Fornecedores que você provavelmente convidaria para um próximo projeto. Citar AAAA – justificativa. Até 200 caracteres, considerando os espaços.}

\subsection{Fornecedores com desempenho abaixo do esperado}
\textcolor{red}{Fornecedores que você não convidaria para um próximo projeto. Citar AAAA – justificativa. Até 200 caracteres, considerando os espaços.}